\uniterablesection{ДОДАТКИ}

\vspace{1em}
\noindent\textbf{Додаток А. Обробник події про виявлену активність}

\vspace{0.5em}
Наведений нижче фрагмент відповідає за створення записів звичок на основі активності, виявленої у зовнішніх сервісах. Він добирає звички користувача за джерелом автоматизації, відсіює вже зафіксовані події за зовнішнім ідентифікатором і створює нові записи (лістинг~\ref{lst:handler}).

\begin{lstlisting}[style=nocodeframe, caption={Створення записів за виявленою активністю}, label={lst:handler}]
public async ValueTask Handle(
    IntegrationActivitiesDetectedIntegrationEvent notification,
    CancellationToken cancellationToken)
{
    var habits = await habitService.GetByAutomationSourceAsync(
        notification.UserId, notification.Source, cancellationToken);

    if (habits.Count == 0)
    {
        return;
    }

    var incomingExternalIds = notification.Activities
        .Select(a => a.ExternalId).ToHashSet();

    var existingExternalIds = await dbContext.HabitEntries.AsNoTracking()
        .Where(e => e.UserId == notification.UserId
            && e.ExternalId != null
            && incomingExternalIds.Contains(e.ExternalId))
        .Select(e => e.ExternalId!)
        .ToHashSetAsync(cancellationToken);

    var newActivities = notification.Activities
        .Where(a => !existingExternalIds.Contains(a.ExternalId))
        .ToList();

    foreach (var habit in habits)
    {
        foreach (var activity in newActivities)
        {
            var result = HabitEntryAggregate.Create(
                Id<HabitEntryAggregate>.NewId(),
                notification.UserId, habit.HabitId,
                value: activity.Value, notes: activity.Notes,
                HabitEntrySource.Automation,
                externalId: activity.ExternalId, activity.OccurredAtUtc);

            if (result.IsFailure)
            {
                continue;
            }

            dbContext.HabitEntries.Add(result.Value.ToEntity());
        }
    }

    await dbContext.SaveChangesAsync(cancellationToken);
}
\end{lstlisting}

\vspace{1em}
\noindent\textbf{Додаток Б. Відбір подій активності GitHub}

\vspace{0.5em}
Цей фрагмент сервісу опитування GitHub демонструє відбір релевантних подій (комітів і запитів на злиття) та публікацію інтеграційної події для подальшої обробки (лістинг~\ref{lst:github}).

\begin{lstlisting}[style=nocodeframe, caption={Відбір подій GitHub та публікація інтеграційної події}, label={lst:github}]
var filteredEvents = eventsResult.Value
    .Where(e =>
        (e.Type == GitHubEventTypes.Push
            || e.Type == GitHubEventTypes.PullRequest)
        && e.CreatedAtUtc > todayDay
        && (e.CreatedAtUtc > lastSyncedAtUtc
            || lastSyncedAtUtc is null))
    .ToList();

var activities = filteredEvents
    .Select(e => new IntegrationActivityItem(e.Id, e.CreatedAtUtc, BuildNotes(e)))
    .ToList();

if (activities.Count == 0)
{
    return Result.Success();
}

await publisher.Publish(
    new IntegrationActivitiesDetectedIntegrationEvent(
        identityUserId, IntegrationActivitySource.Github, activities),
    cancellationToken);

return Result.Success();
\end{lstlisting}

\vspace{1em}
\noindent\textbf{Додаток В. Інструкція користувача}

\vspace{0.5em}
\textit{Передумови.} Для запуску системи потрібні: пакет .NET 10 SDK, середовище Node.js, Docker (для контейнера бази даних, який автоматично піднімає .NET Aspire) та сучасний вебоглядач.

\textit{Запуск серверної частини.} У каталозі оркестратора виконати команду запуску — .NET Aspire самостійно підніме контейнер PostgreSQL та вебсервер API:

\begin{lstlisting}[style=nocodeframe]
cd Dailo.AppHost
dotnet run
\end{lstlisting}

\textit{Запуск клієнтської частини.} У каталозі клієнтського застосунку встановити залежності та запустити сервер розробки; інтерфейс стане доступним за адресою \texttt{http://localhost:4200}:

\begin{lstlisting}[style=nocodeframe]
cd Dailo.ClientApp
npm install
npm start
\end{lstlisting}

\textit{Порядок роботи.}
\begin{enumerate}[noitemsep, topsep=2pt, leftmargin=*, label=\arabic*)]
    \item Зареєструвати обліковий запис або увійти за наявними даними.
    \item У розділі «Звички» створити звички, вказавши тип, періодичність, ціль і, за потреби, джерело автоматизації.
    \item За потреби додати записи вручну в розділі «Записи».
    \item У профілі, на вкладці «Integrations», під'єднати GitHub (увівши персональний токен) та/або Strava (через авторизацію OAuth).
    \item Стежити за прогресом на головній панелі: теплова карта та стрічка активності оновлюються автоматично, зокрема за рахунок автоматично зібраних подій із під'єднаних сервісів.
\end{enumerate}
